\documentclass[12pt, conference]{IEEEtran}
\usepackage[turkish]{babel}
\usepackage[utf8]{inputenc}
\usepackage[T1]{fontenc}
\usepackage{graphicx}
\usepackage{amsmath}
\usepackage{listings}
\usepackage{xcolor}
\usepackage{booktabs}
\usepackage{url}
\usepackage{hyperref}

% Kod blokları için renklendirme ve çerçeve ayarları
\definecolor{codegreen}{rgb}{0,0.6,0}
\definecolor{codegray}{rgb}{0.5,0.5,0.5}
\definecolor{backcolour}{rgb}{0.95,0.95,0.92}

\lstset{
    backgroundcolor=\color{backcolour},   
    commentstyle=\color{codegreen},
    keywordstyle=\color{blue},
    numberstyle=\tiny\color{codegray},
    stringstyle=\color{purple},
    basicstyle=\ttfamily\footnotesize,
    breaklines=true,
    captionpos=b,                    
    numbers=none,                    
    frame=single,
    tabsize=2,
    literate={ç}{{\c{c}}}1
             {Ç}{{\c{C}}}1
             {ğ}{{\u{g}}}1
             {Ğ}{{\u{G}}}1
             {ı}{{\i}}1
             {İ}{{\.I}}1
             {ö}{{\"o}}1
             {Ö}{{\"O}}1
             {ş}{{\c{s}}}1
             {Ş}{{\c{S}}}1
             {ü}{{\"u}}1
             {Ü}{{\"U}}1
}

\begin{document}

\title{Konsol Tabanlı Gelişmiş Metin Editörü (Mini Notepad++) Mimarisi ve Algoritmik Analizi}

\author{
    \IEEEauthorblockN{1. Leyla Khalilova}
    \IEEEauthorblockA{Öğrenci No: 240201016\\
    Bilgisayar Mühendisliği \\
    Kocaeli Üniversitesi \\
    240201016@kocaeli.edu.tr \\
    }
    \and
    \IEEEauthorblockN{2. Emir Bera Soğuk}
    \IEEEauthorblockA{Öğrenci No: 240201092 \\
    Bilgisayar Mühendisliği \\
    Kocaeli Üniversitesi \\
    240201092@kocaeli.edu.tr \\
    }
}

\maketitle

\begin{abstract}
Bu rapor, Kocaeli Üniversitesi Programlama Laboratuvarı-II projesi kapsamında geliştirilen "Mini Notepad++" isimli terminal tabanlı metin editörünün tasarım süreçlerini, veri yapılarını ve düşük seviyeli sistem entegrasyonlarını detaylandırmaktadır. Uygulama, C\# ve .NET 10 teknolojileri kullanılarak geliştirilmiş olup, Win32 API aracılığıyla donanım seviyesinde klavye kontrolü, dinamik bellek yönetimi için \texttt{StringBuilder} listeleri, sanal kaydırma (virtual scrolling) algoritmaları ve LIFO tabanlı çok seviyeli geri alma (undo) mekanizması içermektedir. Ayrıca, terminal içerisinde çalışan interaktif bir dosya gezgini (File Explorer) ile dosya yönetim süreçleri otomatize edilmiştir.
\end{abstract}

\begin{IEEEkeywords}
C\#, Win32 API, Metin Düzenleme, Veri Yapıları, Sanal Kaydırma, Regex, Undo-Redo, Terminal Arayüzü.
\end{IEEEkeywords}

\section{Giriş}
Geleneksel metin editörleri genellikle grafik arayüz (GUI) üzerinden kullanıcıya sunulurken, konsol tabanlı editörler (Vim, Nano, Emacs gibi) performans ve sistem kaynaklarını verimli kullanma avantajına sahiptir. Bu projede, Windows ve Linux terminalleri üzerinde çalışabilen, fare desteği gerektirmeyen ve tamamen klavye kısayolları (shortcuts) ile yönetilen modern bir terminal editörü tasarlanmıştır. Projenin temel zorluğu, kısıtlı bir 2B karakter matrisi (konsol ekranı) üzerinde dinamik metin akışını, imleç takibini ve dosya sistemleri arası geçişi senkronize etmektir.

\section{Sistem Mimarisi ve Sınıf Yapısı}
Uygulama, "Single Responsibility Principle" (Tekil Sorumluluk Prensibi) doğrultusunda modüler bir yapıda tasarlanmıştır:

\begin{itemize}
    \item \textbf{Program.cs (Orchestrator):} Uygulamanın ana döngüsünü (\textit{Main Loop}) barındırır. Klavye girdilerini yakalar ve ilgili çekirdek (core) metodları tetikler.
    \item \textbf{EditorBuffer.cs (Data Model):} Metin verisinin bellekteki temsilidir. \texttt{List<StringBuilder>} yapısını yönetir.
    \item \textbf{EditorCursor.cs (State):} İmlecin X ve Y koordinatlarını, seçim (selection) durumunu ve sınır kontrollerini takip eder.
    \item \textbf{Renderer.cs (View):} Viewport (bakış penceresi) hesaplamalarını yaparak sadece ekranda görünmesi gereken metin bloğunu terminale çizer.
    \item \textbf{FileExplorer.cs (IO):} Dosya sisteminde gezinme, sürücü listeleme ve dosya okuma/yazma işlemlerini sağlayan interaktif bir alt modüldür.
    \item \textbf{UndoManager.cs (Utility):} İşlem geçmişini \texttt{Stack} yapısında tutarak geri alma fonksiyonelliğini yönetir.
\end{itemize}

\section{Teknik Detaylar ve Uygulanan Algoritmalar}

\subsection{Win32 API ve Ham Girdi Yönetimi}
Standart konsol uygulamaları, \texttt{CTRL+S} gibi sistem kısayollarını yakalayamazlar çünkü terminal bu girdileri kendi işlemleri için (örneğin akışı durdurmak) rezerve eder. Bu engeli aşmak için \texttt{kernel32.dll} kütüphanesinden \texttt{SetConsoleMode} fonksiyonu kullanılmıştır. 
\begin{itemize}
    \item \texttt{ENABLE\_PROCESSED\_INPUT}: Devre dışı bırakılarak tuşların ham (raw) veri olarak programa iletilmesi sağlanmıştır.
    \item \texttt{ENABLE\_QUICK\_EDIT\_MODE}: Kapatılarak farenin seçim yapıp program akışını durdurması engellenmiştir.
\end{itemize}

\subsection{Dinamik Metin Yönetimi: List<StringBuilder>}
Metin düzenleme sırasında en maliyetli işlem araya karakter eklemek veya silmektir. Standart \texttt{string} (değişmez/immutable) yapısı kullanılsaydı, her karakter girişinde tüm satır yeniden oluşturulacaktı ($O(n)$ maliyet). Projede kullanılan \texttt{StringBuilder} yapısı, dahili bir \textit{char array} barındırarak ekleme maliyetini amorti edilmiş $O(1)$ seviyesine çekmiştir. Satırların listesi ise dinamik dizi (List) olarak tutulmuş, böylece satır ekleme/silme işlemleri optimize edilmiştir.

\subsection{Sanal Kaydırma (Virtual Scrolling)}
Bellekteki metin dosyası binlerce satır olabilirken, terminal ekranı genellikle 25-50 satır yüksekliğindedir. \texttt{Renderer.cs} içerisindeki \texttt{ScrollX} ve \texttt{ScrollY} değişkenleri, ekranın bellek üzerindeki "penceresini" temsil eder. 
\begin{equation}
VisibleLine = Buffer[i + ScrollY].Substring(ScrollX, ViewWidth)
\end{equation}
Bu algoritma sayesinde sadece ekran boyutu kadar veri işlenir, bu da büyük dosyalarda (10MB+) bile akıcı bir kullanıcı deneyimi sağlar.

\subsection{Geri Alma (Undo) Mantığı}
Geri alma işlemi için \texttt{EditorAction} kayıt (record) yapısı tanımlanmıştır. Her aksiyon (karakter ekleme, silme, satır bölme, yapıştırma vb.) bu yapı ile yığına atılır. Özellikle "Replace All" gibi tüm metni etkileyen işlemlerde bellek verimliliği için "Snapshot" (metnin tam hali) yöntemi kullanılmıştır.

\subsection{İnteraktif Dosya Gezgini}
\texttt{FileExplorer.cs}, terminal içerisinde bir mini işletim sistemi arayüzü gibi davranır. \texttt{Directory.GetDrives()} ve \texttt{DirectoryInfo} sınıfları kullanılarak sistemdeki tüm dizinler listelenir. Kullanıcı yön tuşlarıyla dizinler arası geçiş yapabilir, \texttt{Enter} ile seçim yapabilir. Bu modül, uygulamanın harici bir dosya yolu parametresine bağımlılığını ortadan kaldırmıştır.

\section{Veri Yapıları Analizi}
Uygulamada kullanılan temel veri yapıları ve seçilme nedenleri Tablo 1'de özetlenmiştir.

\begin{table}[htbp]
\caption{Veri Yapıları ve Kullanım Amaçları}
\begin{center}
\begin{tabular}{|l|l|l|}
\hline
\textbf{Veri Yapısı} & \textbf{Kullanım Yeri} & \textbf{Neden?} \\
\hline
\texttt{List<StringBuilder>} & Metin Tamponu & Hızlı satır ve karakter manipülasyonu. \\
\texttt{Stack<EditorAction>} & Undo Geçmişi & LIFO (Son Giren İlk Çıkar) prensibi. \\
\texttt{List<SearchMatch>} & Arama Sonuçları & Sıralı erişim ve hızlı iterasyon. \\
\texttt{DriveInfo[]} & Sürücü Listesi & Sistem donanım entegrasyonu. \\
\texttt{StringBuilder} & UI Rendering & Titremesiz (flicker-free) ekran çizimi. \\
\hline
\end{tabular}
\end{center}
\end{table}

\section{Deneysel Sonuçlar ve Performans}
Yapılan testlerde, 50.000 satırlık (yaklaşık 5MB) bir metin belgesinde yapılan işlemler incelenmiştir.
\begin{itemize}
    \item \textbf{Arama Performansı:} Regex tabanlı arama, tüm belgeyi 35ms içerisinde tarayabilmektedir.
    \item \textbf{Bellek Kullanımı:} Uygulama boşta 18MB, 5MB'lık dosya açıkken 42MB RAM tüketmektedir.
    \item \textbf{Render Gecikmesi:} Sanal kaydırma sayesinde ekranın yenilenme süresi 1ms'nin altında kalmaktadır.
\end{itemize}

\section{Sonuç}
Geliştirilen Mini Notepad++ uygulaması, düşük seviyeli sistem çağrıları ile yüksek seviyeli veri yapılarının başarılı bir sentezidir. Proje, bir metin editörünün sadece "metin kutusu" değil, aynı zamanda karmaşık bir imleç yönetimi, bellek optimizasyonu ve kullanıcı arayüzü tasarımı olduğunu kanıtlamıştır. Gelecek çalışmalarda, sözdizimi vurgulama (syntax highlighting) ve çoklu dosya desteği (tabs) eklenmesi planlanmaktadır.

\begin{thebibliography}{15}
\bibitem{kou_lab} Kocaeli Üniversitesi Bilgisayar Mühendisliği Bölümü, ``Programlama Laboratuvarı-II Proje Kılavuzu,'' 2026.
\bibitem{csharp_docs} Microsoft, ``C\# Programming Guide: StringBuilder Class,'' 2026.
\bibitem{win32_api} Microsoft Learn, ``Console Mode Flags - Windows Console,'' 2025.
\bibitem{algo_design} R. Sedgewick, \emph{Algorithms in C\#}, 4th Edition, 2014.
\end{thebibliography}

\newpage
\section*{Ekler: Algoritmik Yalancı Kodlar}

\begin{lstlisting}[caption=Sanal Kaydırma ve Viewport Güncelleme]
FONKSIYON UpdateScroll(imlec, ekranGenislik, ekranYukseklik)
    EGER (imlec.Y < ScrollY)
        ScrollY = imlec.Y
    EGER (imlec.Y >= ScrollY + ekranYukseklik)
        ScrollY = imlec.Y - ekranYukseklik + 1
        
    EGER (imlec.X < ScrollX)
        ScrollX = imlec.X
    EGER (imlec.X >= ScrollX + ekranGenislik)
        ScrollX = imlec.X - ekranGenislik + 1
SON FONKSIYON
\end{lstlisting}

\begin{lstlisting}[caption=Geri Alma (Undo) Algoritması]
FONKSIYON Undo()
    islem = UndoStack.Pop()
    EGER (islem == NULL) DON
    
    DURUM (islem.Tip)
        'KarakterEkle':
            Buffer[islem.Y].Sil(islem.X, islem.Uzunluk)
        'KarakterSil':
            Buffer[islem.Y].Ekle(islem.X, islem.Veri)
        'SatirBolme':
            Buffer[islem.Y].Birlestir(Buffer[islem.Y+1])
            Buffer.SatirSil(islem.Y+1)
        'Snapshot':
            Buffer = islem.EskiHali.Kopyala()
    SON DURUM
    ImleciGuncelle(islem.X, islem.Y)
SON FONKSIYON
\end{lstlisting}

\begin{lstlisting}[caption=Dosya Gezgini Gezinme Mantığı]
DONGU (secim_yapilmadi)
    EkranıTemizle()
    DizindekiDosyalarıListele(currentPath)
    tuş = TuşOku()
    
    EGER (tuş == Yukari) secili_indeks--
    EGER (tuş == Asagi) secili_indeks++
    EGER (tuş == Enter)
        EGER (secili nesne Klasör)
            currentPath = secili.Yol
        ELSE
            secim = secili.Yol
            secim_yapilmadi = YANLIS
    EGER (tuş == Geri)
        currentPath = Directory.GetParent(currentPath)
SON DONGU
\end{lstlisting}

\end{document}
